La astronomía se ha beneficiado enormemente del crecimiento constante en la capacidad de cómputo y los avances en los campos de la minería de datos y la inteligencia artificial. En particular, las modernas simulaciones hidrodinámicas han proporcionado conjuntos de datos masivos de galaxias con niveles de exactitud sin precedentes.

Aprovechando estos avances, este trabajo presenta los resultados de aplicar técnicas de clustering en el área de descomposición dinámica de galaxias, un campo que busca entender cómo las diferentes partes de una galaxia contribuyen a su movimiento y evolución general. Se comparan los resultados obtenidos con los generados por técnicas ya establecidas en el área. Específicamente, se evaluaron los métodos de Hierarchical Clustering, Fuzzy C-Means y Ensemble Agglomerative Clustering, en conjunto con dos técnicas de eliminación de valores atípicos: una propia de la astronomía (Corte en la mitad del radio de masa) y otra de uso general en minería de datos (Isolation Forest).

Los resultados demuestran que Hierarchical Clustering y Fuzzy C-Means pueden obtener resultados cercanos a los métodos establecidos en el área, aunque con limitaciones significativas, especialmente en la identificación de qué componente corresponde a qué grupo. Por otro lado, se observó que las métricas internas de clustering convencionales no son útiles para analizar el desempeño en este problema debido a la naturaleza particular de los datos astronómicos. En cuanto al Evidence Accumulation Clustering, mostró resultados distantes de los esperados y un alto costo computacional, lo que desaconseja su uso en la forma explorada. Un hallazgo importante fue la incapacidad de determinar si la eliminación de valores atípicos beneficia la tarea de búsqueda de grupos en este contexto específico.

Este trabajo abre camino a futuras investigaciones en el campo, incluyendo la extensión del conjunto de datos para incluir galaxias espirales, la exploración de técnicas de normalización y preprocesamiento adicionales, y la evaluación de otras técnicas de eliminación de valores atípicos y características del conjunto de datos que puedan ser relevantes para la descomposición dinámica de galaxias.

\vfill

\textbf{Palabras Claves:} Minería de datos - Aprendizaje no supervisado - Clustering - Eliminación de atípicos - Simulaciones numéricas - Descomposición dinámica - Parámetro de circularidad